%============================================================
% CH1 · P1 · PROBLEM 1 — FIND THE TOWN JUDGE (LeetCode 997)
%============================================================
\subsection{Problem 1.1.1: Find the Town Judge}
\problemlabel{1.1.1}{ssec:ch1-town-judge}

\problemcard{1.1.1}{Find the Town Judge}{LeetCode 997}{https://leetcode.com/problems/find-the-town-judge/}{\exeasy\ \lc{Easy}}{In-degree and out-degree counting (one net score per vertex)}

\subsubsection*{1. Problem statement}

A town has $n$ people labelled $1..n$. Rumour says one of them is secretly the \textbf{town judge}. If the judge exists, then
\begin{enumerate}[leftmargin=*, label=(\arabic*)]
    \item the judge trusts nobody;
    \item everybody except the judge trusts the judge;
    \item exactly one person satisfies (1) and (2).
\end{enumerate}
You are given the list \texttt{trust}, where \texttt{trust[i]} $= [a_i, b_i]$ means that person $a_i$ trusts person $b_i$ (and no other trust relations exist). Return the label of the judge if the judge exists and can be identified, otherwise $-1$.

\textbf{Constraints:} $1 \le n \le 1000$, \; $0 \le |\texttt{trust}| \le 10^4$, \; all pairs are unique, \; $a_i \ne b_i$, \; $1 \le a_i, b_i \le n$.

\begin{table}[H]
\centering\small
\begin{tabular}{@{}clcl@{}}
\toprule
$n$ & \textbf{trust} & \textbf{Output} & \textbf{Why} \\
\midrule
2 & $[[1,2]]$ & 2 & 1 trusts 2, 2 trusts nobody \\
3 & $[[1,3],[2,3]]$ & 3 & both others trust 3 \\
3 & $[[1,3],[2,3],[3,1]]$ & $-1$ & 3 is trusted by all, but trusts 1 \\
\bottomrule
\end{tabular}
\end{table}

\begin{diagrambox}
\begin{tikzpicture}[x=1cm, y=1cm]
    \begin{scope}
        \node[gv] (a1) at (0,0) {1}; \node[gvdst] (a2) at (1.6,0) {2};
        \draw[gde] (a1) -- (a2);
        \node[font=\scriptsize\bfseries] at (0.8,-1.45) {Example 1: judge 2};
    \end{scope}
    \begin{scope}[xshift=4.2cm]
        \node[gv] (b1) at (0,0.6) {1}; \node[gv] (b2) at (0,-0.6) {2}; \node[gvdst] (b3) at (1.6,0) {3};
        \draw[gde] (b1) -- (b3); \draw[gde] (b2) -- (b3);
        \node[font=\scriptsize\bfseries] at (0.8,-1.45) {Example 2: judge 3};
    \end{scope}
    \begin{scope}[xshift=8.6cm]
        \node[gv] (c1) at (0,0.6) {1}; \node[gv] (c2) at (0,-0.6) {2}; \node[gvdead] (c3) at (1.6,0) {3};
        \draw[gde] (c1) to[bend left=15] (c3); \draw[gde] (c2) -- (c3);
        \draw[gback] (c3) to[bend left=15] (c1);
        \node[font=\scriptsize\bfseries] at (0.8,-1.45) {Example 3: no judge ($-1$)};
    \end{scope}
\end{tikzpicture}
\figcap{The three official examples. An arrow $a \to b$ means ``$a$ trusts $b$''. In Example 3 the dashed arrow $3 \to 1$ disqualifies 3.}
\end{diagrambox}

\subsubsection*{2. The graph model}
\begin{modelbox}[Graph Model --- Find the Town Judge]
\renewcommand{\arraystretch}{1.35}
\begin{tabularx}{\linewidth}{@{}>{\bfseries}l X@{}}
Vertices & the $n$ people, labelled $1..n$ \\
Edges & a \emph{directed} edge $a \to b$ for every pair $[a,b]$ in \texttt{trust}; stored only as the given edge list \\
Weights & none \\
Question & find the vertex $v$ with $\indeg{v} = n-1$ and $\outdeg{v} = 0$ (a sink that every other vertex points to), or report that none exists \\
Size & $|V| \le 1000$, $|E| \le 10^4$: anything up to $\approx 10^7$ operations is fine, $\bigO{|V|+|E|}$ is trivial \\
\end{tabularx}
\end{modelbox}

\begin{diagrambox}
\begin{tikzpicture}[x=1cm, y=1cm]
    \begin{scope}
        \node[font=\small\bfseries, text=deepblue] at (0,1.6) {the rules};
        \node[tag, text width=4.6cm, align=left] at (0,0.2)
            {(1) the judge trusts nobody\\(2) everybody else trusts the judge};
        \draw[flow] (2.6,0.2) -- node[lbl, above]{translate} (4.4,0.2);
    \end{scope}
    \begin{scope}[xshift=7.4cm]
        \node[gvdst] (J) at (0,0) {$J$};
        \foreach \a [count=\i] in {90,150,210,270,330,30}{
            \node[gv, minimum size=5.5mm] (p\i) at (\a:1.45) {};
            \draw[gde] (p\i) -- (J);
        }
        \node[gdist, anchor=west] at (1.9,0.35) {$\indeg{J} = n-1$};
        \node[gdist, anchor=west] at (1.9,-0.25) {$\outdeg{J} = 0$};
    \end{scope}
\end{tikzpicture}
\figcap{The judge is a vertex with every other vertex pointing to it and no arrow leaving it. What the other people trust among themselves does not matter.}
\end{diagrambox}

\subsubsection*{3. Key observations}

\begin{itemize}[leftmargin=*]
    \item \textbf{Both rules are degree conditions.} ``Trusts nobody'' is $\outdeg{v} = 0$; ``trusted by everybody else'' is $\indeg{v} = n-1$, because the pairs are unique and nobody trusts himself, so the $n-1$ arrows into $v$ must come from $n-1$ different people. No path, no traversal --- Pattern 1.
    \item \textbf{Rule (3) is automatic.} Suppose $J_1 \ne J_2$ both satisfy (1) and (2). By (2) for $J_1$, $J_2$ trusts $J_1$ --- but by (1) for $J_2$, $J_2$ trusts nobody. Contradiction. So at most one vertex can pass the degree test, and the first one we find is the answer.
    \item \textbf{In-degree alone is not enough.} In Example 3 vertex 3 has $\indeg{3} = 2 = n-1$ but trusts 1. Both conditions must be checked.
    \item \textbf{$n = 1$, no edges:} the single person trusts nobody and is trusted by all $0$ others --- he \emph{is} the judge; answer 1.
\end{itemize}

\begin{diagrambox}
\begin{tikzpicture}[x=1cm, y=1cm]
    \node[gvdst] (J1) at (0,0) {$J_1$};
    \node[gvdst] (J2) at (3,0) {$J_2$};
    \draw[gde, draw=warningred, line width=1.3pt] (J2) to[bend right=25] node[lbl, above]{required by (2) for $J_1$} (J1);
    \node[note, text width=4.2cm, anchor=west] at (4,0) {but (1) for $J_2$ forbids any arrow leaving $J_2$ \;\Large\contradiction};
\end{tikzpicture}
\figcap{Why there can never be two judges: the arrow $J_2 \to J_1$ would have to exist and not exist at the same time.}
\end{diagrambox}

\subsubsection*{4. Approaches}

\begin{bruteforcebox}
\textbf{Idea.} Try every person $c$ as the candidate; scan the whole trust list to check that $c$ trusts nobody and count how many people trust $c$.
\begin{lstlisting}[style=cpp]
#include <bits/stdc++.h>
using namespace std;

class Solution {
public:
    int findJudge(int n, vector<vector<int>>& trust) {
        for (int c = 1; c <= n; c++) {                 // candidate judge
            bool trustsSomeone = false;
            int trustedBy = 0;
            for (auto& t : trust) {
                if (t[0] == c) trustsSomeone = true;   // rule (1) broken
                if (t[1] == c) trustedBy++;
            }
            if (!trustsSomeone && trustedBy == n - 1) return c;
        }
        return -1;
    }
};
\end{lstlisting}
\begin{complexitybox}
\textbf{Time:} $\bigO{n \cdot |E|}$ --- up to $1000 \times 10^4 = 10^7$ steps: still \ac{} here, but it rescans the same list $n$ times to compute numbers that one scan already contains. \textbf{Space:} $\bigO{1}$.
\end{complexitybox}
\end{bruteforcebox}

\begin{algorithmbox}
\textbf{Idea.} Both counts come from one pass. Even better, combine them into one \textbf{net score} $\texttt{score}[v] = \indeg{v} - \outdeg{v}$: every pair $[a,b]$ does \texttt{score[a]--} and \texttt{score[b]++}. The judge is the vertex with score exactly $n-1$.
\begin{lstlisting}[style=cpp]
#include <bits/stdc++.h>
using namespace std;

class Solution {
public:
    int findJudge(int n, vector<vector<int>>& trust) {
        vector<int> score(n + 1, 0);       // score[v] = in(v) - out(v), labels 1..n
        for (auto& t : trust) {
            score[t[0]]--;                 // t[0] trusts someone: out-degree + 1
            score[t[1]]++;                 // t[1] is trusted:     in-degree  + 1
        }
        for (int v = 1; v <= n; v++)
            if (score[v] == n - 1) return v;
        return -1;
    }
};
\end{lstlisting}
\begin{complexitybox}
\textbf{Time:} $\bigO{n + |E|}$ --- one pass over the pairs, one over the people. \textbf{Space:} $\bigO{n}$ for \texttt{score}. \textbf{Verdict:} \ac{}.
\end{complexitybox}
\end{algorithmbox}

\begin{invariantbox}[Why score $= n-1$ identifies the judge]
After the loop, $\texttt{score}[v] = \indeg{v} - \outdeg{v}$ (each pair changes exactly the two counters of its endpoints, see the invariant of Pattern 1). Since the pairs are unique and $a_i \ne b_i$, at most $n-1$ different people can trust $v$, so $\indeg{v} \le n-1$; and $\outdeg{v} \ge 0$. Hence
\[
\texttt{score}[v] = n-1 \iff \indeg{v} = n-1 \text{ and } \outdeg{v} = 0,
\]
which is exactly rules (1) and (2). By the observation ``at most one judge'', the first such $v$ is the unique answer. \hfill$\blacksquare$
\end{invariantbox}

\begin{alternativebox}
\textbf{Two separate counters} (\texttt{in[]}, \texttt{out[]}, test \texttt{in[v] == n-1 \&\& out[v] == 0}) is the same algorithm, slightly longer but robust when duplicates are allowed (then \texttt{score} could reach $n-1$ by accident).

\textbf{Candidate elimination (the ``celebrity'' algorithm).} Store trust as an $n \times n$ boolean matrix. Walk once over the people, keeping one candidate. Comparing the candidate $c$ with the next person $v$ always eliminates one of them:
\begin{itemize}[leftmargin=*]
    \item if $c$ trusts $v$, then $c$ is not the judge (the judge trusts nobody) $\Rightarrow$ the new candidate is $v$;
    \item otherwise $v$ is not the judge (the judge is trusted by everybody, including $c$) $\Rightarrow$ keep $c$.
\end{itemize}
After the walk only one person can still be the judge; check him explicitly.
\begin{lstlisting}[style=cpp]
#include <bits/stdc++.h>
using namespace std;

class Solution {
public:
    int findJudge(int n, vector<vector<int>>& trust) {
        vector<vector<char>> T(n + 1, vector<char>(n + 1, 0));   // T[a][b] = a trusts b
        for (auto& t : trust) T[t[0]][t[1]] = 1;
        int cand = 1;
        for (int v = 2; v <= n; v++)
            if (T[cand][v]) cand = v;              // cand trusts v: cand eliminated
        for (int v = 1; v <= n; v++) {             // verify the survivor
            if (v == cand) continue;
            if (T[cand][v] || !T[v][cand]) return -1;
        }
        return cand;
    }
};
\end{lstlisting}
This uses $\bigO{n^2}$ memory and time here, so it is worse for this problem. Its point: it asks only $3(n-1)$ questions ``does $a$ trust $b$?'' --- the right tool when trust is hidden behind an expensive query function instead of being listed (variant V3).

\begin{diagrambox}
\begin{tikzpicture}[x=1.55cm, y=1cm]
    \node[gvdead] (p1) at (0,0) {1}; \node[gvdead] (p2) at (1.6,0) {2};
    \node[gvdst] (p3) at (3.2,0) {3}; \node[gvdead] (p4) at (4.8,0) {4};
    \node[lbl, text=darkgray] at (0.8,0.75) {1 vs 2};
    \node[lbl, text=darkgray] at (2.4,0.75) {1 vs 3};
    \node[lbl, text=darkgray] at (4.0,0.75) {3 vs 4};
    \node[note, text width=2.1cm] at (0.8,-0.95) {1 does not trust 2: drop 2};
    \node[note, text width=2.1cm] at (2.4,-0.95) {1 trusts 3: drop 1};
    \node[note, text width=2.1cm] at (4.0,-0.95) {3 does not trust 4: drop 4};
    \node[note, anchor=west, text width=2.6cm, align=left] at (5.3,0) {survivor 3: verified, it is the judge};
\end{tikzpicture}
\figcap{Elimination on the trace example below ($n=4$, pairs $1{\to}3, 1{\to}4, 2{\to}3, 2{\to}4, 4{\to}3$): three comparisons leave one candidate.}
\end{diagrambox}
\end{alternativebox}

\begin{tricksbox}
\begin{itemize}[leftmargin=*]
    \item \textbf{Early exit:} if $|\texttt{trust}| < n-1$, nobody can be trusted by $n-1$ people --- return $-1$ immediately (careful: $n = 1$ with no pairs must still return 1, and indeed $0 < 0$ is false).
    \item \textbf{One array instead of two} halves the memory and makes the final test a single comparison --- valid only because pairs are unique.
    \item The \texttt{vector<vector<int>>\&} input is read-only here; iterate with \texttt{auto\&} to avoid copying each pair.
\end{itemize}
\end{tricksbox}

\subsubsection*{5. Trace}

\begin{tracebox}[{Net scores for $n = 4$, trust $= [1{\to}3,\ 1{\to}4,\ 2{\to}3,\ 2{\to}4,\ 4{\to}3]$}]
\centering
\begin{tikzpicture}[x=1cm, y=1cm]
    % snapshot A: after 2 edges
    \begin{scope}
        \node[gvseen] (a1) at (0,1) {1}; \node[gv] (a2) at (0,-0.4) {2};
        \node[gvcur] (a3) at (1.7,1) {3}; \node[gvcur] (a4) at (1.7,-0.4) {4};
        \draw[gde] (a1) -- (a3); \draw[gde] (a1) -- (a4);
        \node[gdist, left=1mm of a1] {$-2$}; \node[gdist, left=1mm of a2] {$0$};
        \node[gdist, right=1mm of a3] {$1$}; \node[gdist, right=1mm of a4] {$1$};
        \node[font=\scriptsize\bfseries] at (0.85,-1.25) {after $1{\to}3,\ 1{\to}4$};
    \end{scope}
    % snapshot B: after 4 edges
    \begin{scope}[xshift=4.3cm]
        \node[gvseen] (b1) at (0,1) {1}; \node[gvseen] (b2) at (0,-0.4) {2};
        \node[gvcur] (b3) at (1.7,1) {3}; \node[gvcur] (b4) at (1.7,-0.4) {4};
        \draw[gde] (b1) -- (b3); \draw[gde] (b1) -- (b4); \draw[gde] (b2) -- (b3); \draw[gde] (b2) -- (b4);
        \node[gdist, left=1mm of b1] {$-2$}; \node[gdist, left=1mm of b2] {$-2$};
        \node[gdist, right=1mm of b3] {$2$}; \node[gdist, right=1mm of b4] {$2$};
        \node[font=\scriptsize\bfseries] at (0.85,-1.25) {after $2{\to}3,\ 2{\to}4$: a tie};
    \end{scope}
    % snapshot C: final
    \begin{scope}[xshift=8.6cm]
        \node[gvdead] (c1) at (0,1) {1}; \node[gvdead] (c2) at (0,-0.4) {2};
        \node[gvdst] (c3) at (1.7,1) {3}; \node[gvdead] (c4) at (1.7,-0.4) {4};
        \draw[gde] (c1) -- (c3); \draw[gde] (c1) -- (c4); \draw[gde] (c2) -- (c3); \draw[gde] (c2) -- (c4);
        \draw[gpath, -{Stealth[length=2.6mm]}] (c4) -- (c3);
        \node[gdist, left=1mm of c1] {$-2$}; \node[gdist, left=1mm of c2] {$-2$};
        \node[gdist, right=1mm of c3] {$\mathbf{3}$}; \node[gdist, right=1mm of c4] {$1$};
        \node[font=\scriptsize\bfseries] at (0.85,-1.25) {after $4{\to}3$: judge 3};
    \end{scope}
\end{tikzpicture}
\figcap{Green numbers are the scores. The last pair breaks the tie: 4 now trusts somebody (score $2 \to 1$) and 3 reaches $n-1 = 3$.}

\medskip
\small
\begin{tabular}{@{}lcccc@{}}
\toprule
\textbf{pair processed} & \texttt{score[1]} & \texttt{score[2]} & \texttt{score[3]} & \texttt{score[4]} \\
\midrule
(start) & 0 & 0 & 0 & 0 \\
$1 \to 3$ & $-1$ & 0 & 1 & 0 \\
$1 \to 4$ & $-2$ & 0 & 1 & 1 \\
$2 \to 3$ & $-2$ & $-1$ & 2 & 1 \\
$2 \to 4$ & $-2$ & $-2$ & 2 & 2 \\
$4 \to 3$ & $-2$ & $-2$ & \textbf{3} & 1 \\
\bottomrule
\end{tabular}
\end{tracebox}

\begin{tracebox}[{Example 3: $n = 3$, trust $= [1{\to}3,\ 2{\to}3,\ 3{\to}1]$}]
\centering
\begin{tikzpicture}[x=1cm, y=1cm]
    \begin{scope}
        \node[gvseen] (a1) at (0,0.7) {1}; \node[gvseen] (a2) at (0,-0.7) {2}; \node[gvcur] (a3) at (1.8,0) {3};
        \draw[gde] (a1) -- (a3); \draw[gde] (a2) -- (a3);
        \node[gdist, left=1mm of a1] {$-1$}; \node[gdist, left=1mm of a2] {$-1$}; \node[gdist, right=1mm of a3] {$2$};
        \node[font=\scriptsize\bfseries] at (0.9,-1.45) {after two pairs: 3 has score $2 = n-1$};
    \end{scope}
    \draw[flow] (3.3,0) -- node[lbl, above]{$3 \to 1$} (5,0);
    \begin{scope}[xshift=6.6cm]
        \node[gv] (b1) at (0,0.7) {1}; \node[gvdead] (b2) at (0,-0.7) {2}; \node[gvdead] (b3) at (1.8,0) {3};
        \draw[gde] (b1) to[bend left=12] (b3); \draw[gde] (b2) -- (b3);
        \draw[gback] (b3) to[bend left=12] (b1);
        \node[gdist, left=1mm of b1] {$0$}; \node[gdist, left=1mm of b2] {$-1$}; \node[gdist, right=1mm of b3] {$1$};
        \node[font=\scriptsize\bfseries] at (0.9,-1.45) {final: no score equals 2 $\Rightarrow -1$};
    \end{scope}
\end{tikzpicture}
\figcap{One outgoing arrow costs the candidate a point: the net score automatically enforces rule (1).}
\end{tracebox}

\subsubsection*{6. Test suite}

\begin{testbox}
\centering\small
\begin{tabular}{@{}cllc@{}}
\toprule
$n$ & \textbf{trust} & \textbf{Why it is interesting} & \textbf{Expected} \\
\midrule
1 & $[\,]$ & single person, no pairs: he is the judge & 1 \\
2 & $[\,]$ & nobody trusts anybody & $-1$ \\
2 & $[[1,2]]$ & official Example 1 & 2 \\
3 & $[[1,3],[2,3],[3,1]]$ & in-degree is $n-1$, but out-degree is 1 & $-1$ \\
3 & $[[1,2],[2,3]]$ & a chain: 3 is trusted by one person only & $-1$ \\
4 & $[[1,2],[2,3],[3,4],[4,1]]$ & a cycle: every score is 0 & $-1$ \\
4 & $[[1,3],[2,3],[4,3],[1,2]]$ & extra trust between non-judges is fine & 3 \\
1000 & $i \to 1000$ for all $i < 1000$ & largest $n$, judge has the last label & 1000 \\
\bottomrule
\end{tabular}
\end{testbox}

\subsubsection*{7. Common mistakes}

\begin{warningbox}[Mistake 1 --- checking only the in-degree]
\begin{lstlisting}[style=snippet]
if (in[v] == n - 1) return v;                    // WRONG: Example 3 returns 3
if (in[v] == n - 1 && out[v] == 0) return v;     // RIGHT
\end{lstlisting}
\end{warningbox}

\begin{warningbox}[Mistake 2 --- special-casing the empty list]
\texttt{if (trust.empty()) return -1;} fails for $n = 1$, where the answer is 1. The general loop handles it: $\texttt{score}[1] = 0 = n - 1$.
\end{warningbox}

\begin{warningbox}[Mistake 3 --- an array of size $n$]
Labels are $1..n$. \texttt{vector<int> score(n)} makes \texttt{score[n]} an out-of-bounds write. Use \texttt{n + 1} slots.
\end{warningbox}

\begin{warningbox}[Mistake 4 --- trusting the net score when duplicates are possible]
If the same pair could appear twice, a person trusted twice by one friend could reach score $n-1$ without being trusted by everybody. LeetCode promises unique pairs; in other judges deduplicate first or count distinct trusters.
\end{warningbox}

\subsubsection*{8. Takeaways}

\begin{itemize}[leftmargin=*]
    \item Translate every rule into a statement about \emph{one vertex and its edges}: then degrees answer it in $\bigO{|V|+|E|}$.
    \item ``Points to nobody'' $=$ out-degree 0, ``everybody points to it'' $=$ in-degree $n-1$ (in a simple directed graph).
    \item Two counters can often be merged into one net score when one of them is bounded.
    \item Uniqueness claims (``exactly one person'') are often consequences of the rules --- prove them, then stop at the first hit.
\end{itemize}

\begin{variantsbox}
\variant{Find Center of Star Graph}{\varin{https://leetcode.com/problems/find-center-of-star-graph/}}
  {undirected: the edges form a star, find the vertex adjacent to all others.}
  {model: the center is the unique vertex of degree $n-1$; it appears in \emph{every} edge.}
  {algorithm: degree counting works; even simpler, return the common endpoint of the first two edges ($\bigO{1}$).}
\variant{Vertices that reach everything}{\varin{https://leetcode.com/problems/minimum-number-of-vertices-to-reach-all-nodes/}}
  {in a DAG, find the smallest set of vertices from which all vertices are reachable.}
  {model: a vertex with in-degree 0 can be reached from nobody, so it must be chosen; every other vertex is reached from some source.}
  {algorithm: count in-degrees, output all $v$ with $\indeg{v} = 0$ --- the same counting loop as the judge.}
\variant{Celebrity with a query oracle}{\varnone}
  {trust is hidden; you may only call \texttt{knows(a,b)} and must use $\bigO{n}$ calls.}
  {model: the same directed graph, but edges are revealed one query at a time.}
  {algorithm: candidate elimination ($n-1$ queries) plus verification ($2(n-1)$ queries) --- Approach 3.}
\variant{Duplicated trust records}{\varnone}
  {the same pair $[a,b]$ may be listed several times.}
  {model: a directed multigraph; the judge condition is about \emph{distinct} trusters.}
  {algorithm: sort and deduplicate the pairs (or use a set) first, then count; the net score is unsafe on the raw list.}
\end{variantsbox}
